\chapter{Fix 47: The Finite Volume Luscher Correction}
\label{ch:fix47_finite_volume}

\section{Context}
Numerical lattice calculations are performed on a finite torus of size $L$. To relate these to the infinite volume mass gap $\Delta$, we must prove the convergence is exponential, ruling out massless modes that would cause polynomial corrections.

\section{Theorem O.1 (Exponential Finite Volume Scaling)}
If the infinite volume theory has a mass gap $m$, the mass gap $m(L)$ on a torus of size $L$ satisfies:
\begin{equation}
m(L) = m - \frac{A}{L} e^{-m L} (1 + O(e^{-m L}))
\end{equation}
This validates the extraction of the physical gap from finite lattice sequences.

\section{Proof Derivation}
\subsection{1. Self-Energy on the Cylinder}
We model the finite volume effect as a "virtual particle" traveling around the world-circle of length $L$. The shift in mass is the self-energy of the flux tube due to this wrapping.

\subsection{2. Scattering Amplitude Representation}
The mass shift is determined by the forward scattering amplitude $F(\theta)$ of the particle with itself:
\begin{equation}
\delta m \propto \int F(i\theta) e^{-m L \cosh \theta} d\theta
\end{equation}.

\subsection{3. Saddle Point Evaluation}
The integral is dominated by the saddle point at $\theta=0$. The exponential factor $e^{-mL}$ controls the magnitude.

\subsection{4. Bethe-Salpeter Compactness}
Using the compactness of the Bethe-Salpeter kernel established in \textbf{Appendix AU}, we ensure $F$ is analytic in the relevant strip. This proves that the finite volume corrections are exponentially suppressed, confirming the gap is a bulk property of the vacuum.
